% ============================================================
\section{Discussion}
\label{sec:discussion}
% ============================================================

\subsection{Sources of Edge and Economic Mechanism}

The strategy is built for a drawdown-constrained investor, one whose
objective is capital preservation while still participating in trends, not
unconditional Sharpe maximisation. The nested attribution
(Table~\ref{tab:ladder}) decomposes the edge accordingly. The bulk of return
comes from momentum, a well-documented factor we do not claim as a contribution.
Concentrating to three assets and imposing the absolute threshold each add
economically meaningful but individually marginal increments, and minimum-correlation
selection adds little return while compressing risk. No single
rule is statistically significant on its own out-of-sample (momentum: $p = 0.80$,
min-corr isolated: $p = 0.68$ on the locked test); the edge emerges from
stacking complementary rules (total: $p = 0.008$ on the test block). That
is both a more honest description than pinning everything on one mechanism and a
more defensible one against the charge that a single tuned parameter drives the result.

Two economic mechanisms drive the non-momentum component. The first is the breadth of
momentum, which works as an endogenous regime indicator. Momentum reflects persistent,
slow-moving demand for an asset class; when it collapses simultaneously
across a diversified cross-section (fewer than three assets clearing the absolute
threshold), it signals a broad risk-off transition rather than idiosyncratic
rotation. Rotating to duration or cash in that state is the flight-to-quality
response \citep{barroso2015, daniel2016}, taken with the same signal that
drives offensive allocation. No separate, separately-tunable market-timing rule
is needed, which keeps the strategy simple and auditable. The second mechanism is
minimum-correlation selection, which hedges against the well-documented rise in
cross-asset correlations during stress. Because rank within the eligible set carries
no further expected-return information (Empirical Regularity~\ref{lem:homogeneity}),
the only economically relevant lever left is the covariance structure. Choosing the
lowest-correlation triplet ex ante lowers variance and thins the left tail without
altering expected return, the risk-control (not return) pattern the data show.

The strategy is also not merely a repackaging of the trend premium. When we
control for a time-series-momentum factor built from the same 15 assets
(Table~\ref{tab:factors}), a residual alpha of roughly $+4$--$8\%$ per annum
survives, regardless of whether the TSMOM factor is specified as long-short
($+4.3\%$/yr, $p = 0.04$ on the locked test) or, more rigorously, as long-only
(the appropriate specification for a long-only strategy: $+7.6\%$/yr, $p = 0.018$
on the locked test). Controlling for the long-only TSMOM factor gives a
larger residual alpha than the long-short version, because the long-only
factor absorbs more of GAT's variance ($R^2 \approx 0.52$) while leaving a larger
return component unexplained. The non-trend component is the combination of
concentration, the breadth-based defensive switch, and minimum-correlation
construction, which together convert a known return source into a drawdown-controlled
allocation that is not fully spanned by any passive trend strategy on the same
assets.

\subsection{Honest Benchmarking}

The $1/N$ portfolio is a low bar by design. \citet{demiguel2009} adopt it
because it is hard to beat net of estimation error, not because it is a
strong risk-adjusted competitor, so we do not rest the case on it. On the locked
test the strategy out-performs $1/N$ only marginally on the Sharpe
difference ($p = 0.079$, a 60-month window with limited power) but robustly on the
spanning alpha ($p = 0.008$). The more demanding evidence is the alpha against
recognised factor models: it survives the CAPM, a two-factor equity-plus-duration
model, and a time-series-momentum factor built from the same assets
(Tables~\ref{tab:alpha} and~\ref{tab:factors}).
It also beats a Faber (2007) trend rule ($p = 0.032$), so the result is not pure
trend-following. It does \emph{not} beat 60/40 or the S\&P~500 on Sharpe, where
60/40 benefits from the duration tailwind of the 2010s. Against
these benchmarks the case is drawdown, not Sharpe: comparable or higher
risk-adjusted return at a materially lower maximum drawdown. That advantage is
statistically significant against the S\&P~500 and $1/N$ and marginal
against 60/40 ($\Delta$MaxDD $+17.7$pp vs SPY, $p=0.018$; $+16.7$pp vs $1/N$,
$p=0.039$; $+14.0$pp vs 60/40, $p=0.056$; locked test block), given the high sampling
variability of maximum drawdown and 60/40's own bond cushion.
We state this explicitly rather than select the benchmark that flatters the
strategy. On portfolio construction, the equal-weight minimum-correlation rule
captures the diversification benefit of optimised construction while avoiding the
estimation error of covariance-matrix inversion \citep{demiguel2009}; none of the
optimised alternatives (minimum-variance, risk-parity, inverse-volatility) adds a
spanning alpha over GAT or achieves a lower drawdown at comparable Sharpe
(Section~\ref{sec:res_construction}).

\subsection{Data Snooping and Researcher Degrees of Freedom}

Four parameter dimensions enter the design, so the apparent performance of the
selected configuration could in principle reflect a search over a large grid
\citep{white2000, harvey2018}. We address this on two independent fronts.
First, the way each parameter is set leaves little room to overfit
(Section~\ref{sec:methodology}). Three of the four are derived, and their
training optimum coincides with both the ex-ante literature value and the pooled
train+validation optimum (the momentum signal, the Top-$N$ value, and the triplet
size), so they carry no free degree of freedom. The fourth, the absolute threshold,
is the one most exposed to a snooping critique and we therefore do \emph{not}
optimise it at all: it is fixed ex ante at the risk-free rate ($\theta=2\%$). We are
candid that the in-sample objective would actually prefer a higher threshold
($\approx 4$--$5\%$); declining to chase that maximum is a deliberate restriction of
researcher degrees of freedom, not a coincidental optimum. Second, the
matched-permutation attribution (Table~\ref{tab:attribution}) shows that only
momentum and minimum correlation beat a random matched baseline. The two rules that
drive the result are exactly the two with an independent economic rationale, while
the threshold contributes through risk control rather than selection
(Section~\ref{sec:theta}); the parameters most exposed to a snooping critique are
the ones that are either inert or fixed ex ante.

We also declare, and then close, the limitations of the data construction. The
proxy calibration uses the full overlap and therefore embeds a bounded look-ahead
in the early training segment (2004--2007). Three results bound its impact
(Section~\ref{sec:res_robustness}): the validation and test blocks are 100\% real
ETF data; dropping BWX and RWX (the two assets that previously relied on weak
proxies, now reconstructed from an exact index and a long-history fund) leaves the
alpha at $+5.57\%$/yr ($p = 0.0066$); and, most directly, running the strategy on 100\%
real ETF prices with no proxies at all (2009--2025) yields essentially the same
out-of-sample alpha ($+6.04\%$/yr, $p = 0.0001$). The proxy extension is thus a
power-enhancing device, not the source of the edge. An expanding-window
recalibration would close the residual concern formally, but the real-data result
already shows it does not change the conclusion.

The \citet{bailey2014} Deflated Sharpe Ratio at the $N=24$ searched configurations
(the threshold is fixed ex ante and adds none) is $0.975$, above the $0.95$
threshold, and remains so for any assumed trial count up to $N \approx 154$; the
in-sample Sharpe therefore survives the data-snooping adjustment. We nonetheless
treat the out-of-sample alpha and matched-permutation attribution as the primary
defence against data snooping.

A transparency note on pre-registration. We did not formally pre-register the
strategy before running the test block. The decision sequence is documented in
Section~\ref{sec:theta}: parameters were locked on the training grid and the absolute
threshold was fixed ex ante at $\theta = 2\%$ without examining test returns; the
test block was then evaluated once. We acknowledge that, in the absence of a
time-stamped pre-registration, this sequential process cannot be externally verified
and therefore declare it as a stated procedural constraint, not a resolved concern.

A second transparency note concerns the B2 intermediate model. B2 (top-3 momentum
of the eligible set, with defensive switch, but without minimum-correlation
selection) generates $\alpha = +9.98\%$/yr versus $1/N$ on the locked test
($p = 0.0002$) and a Sharpe of $1.151$ versus GAT's $0.873$. This finding is not
post-hoc: B2 is an intermediate step in the nested ladder, not a separately
discovered alternative, and the final model GAT was determined before the test block
was examined. Nonetheless, B2 outperforms GAT on the Sharpe criterion on this
particular test window. The explanation is precisely the stated design objective:
minimum-correlation selection sacrifices Sharpe for drawdown compression
(B2: $-8.30\%$, GAT: $-7.06\%$) and correlation reduction ($\bar\rho$: 0.50 vs 0.22).
On the longer post-training window (2016--2025) B2 leads on Sharpe by a smaller margin
and the incremental step GAT$\sim$B2 is alpha-neutral in both directions
($+0.23\%$ OOS, $-0.82\%$ test). We report this candidly: the minimum-correlation rule
achieves its design purpose (risk control, not return) but a researcher selecting on
the test Sharpe alone would choose B2 over GAT. We do not make that claim, and we
have not adjusted the strategy ex post.

\subsection{Limitations}

Several limitations bound the interpretation of these results.

\textit{Backtest overfitting.}
All parameters were selected within the training window. The strict temporal split
and the attribution tests mitigate the most direct forms of look-ahead and
snooping, but the researcher degrees of freedom in choosing the overall
architecture are not fully captured by any in-sample control. A prospective live
track record is the only complete remedy.

\textit{Data construction.}
The extended history relies on proxies; every proxied series now uses its exact
underlying index or a source with $R^2 \geq 0.92$, the single exception being the
2003--2006 segment of the mortgage-REIT series (REM, $R^2 = 0.80$), which rarely
enters the triplet; the 11-month pre-launch segment of TIP uses a real-yield
constant-maturity series ($R^2 = 0.40$, a short segment that precedes the strategy's
live window). The calibration look-ahead is bounded and, as shown above, does not
drive the result (the real-data alpha is unchanged); an expanding-window
recalibration would close the concern formally and is left for future work.

\textit{Momentum homogeneity.}
Lemma~\ref{lem:homogeneity} holds in training and validation but shows a weak
negative rank effect in the locked test window (Spearman $\hat{\rho} = -0.14$,
$p = 0.010$; ANOVA $p = 0.25$). The negative sign is inconsistent with momentum
within the eligible set and is directionally favourable to minimum-correlation
selection over rank-based selection, but the result weakens the claim that the
homogeneity premise is an out-of-sample regularity rather than a training-period
property. A prospective sample would provide a cleaner test.

\textit{Single universe and turnover.}
The strategy uses a fixed 15-asset universe and does not adapt as new instruments
become liquid. Its turnover ($9.83\times$/yr) is moderate-to-high; the no-trade-band
variant reduces it to $7.66\times$/yr while preserving the alpha, but the strategy is
inherently more cost-sensitive than low-turnover peers. The audited break-even cost
is near $62$ bps one-way (roughly six times the realistic figure for liquid
ETFs); the edge is robust to realistic costs but loses significance at a conservative
$35$ bps. Calendar (quarterly) rebalancing is not a viable reduction because it
sharply worsens the drawdown.

\textit{Regime dependence.}
The out-of-sample advantage is concentrated in stress and sideways regimes; in
sustained bond or equity bull markets the strategy merely keeps pace with passive
benchmarks while protecting against drawdowns that those benchmarks do not.

\textit{Minimum-correlation rule: design vs.\ realised test performance.}
The minimum-correlation selection rule adds a positive but insignificant return
increment over the longer post-training window ($+0.23\%$/yr, $p = 0.869$) and a
small negative increment on the 60-month locked test ($-0.82\%$, $p = 0.682$). Its
value is risk reduction, not return, which is exactly what Proposition~\ref{prop:mincorr}
predicts and what Section~\ref{sec:res_mincorr} documents. The implication for
practitioners is that a return-focused investor should use B2 (top-3 momentum of the
eligible set) rather than GAT; GAT is the appropriate choice for a
drawdown-constrained investor. We designed GAT for the latter objective and the test
block confirms the design works as specified, but we do not claim the minimum-correlation
rule improves risk-adjusted return.
